\section{Competence: is it actually good at this, or only good on paper?}
\label{sec:competence}

The first question a clinician asks of any tool is whether it is competent at the
specific task, on patients like the ones in this clinic. Competence is not a single
benchmark number; it is a stack of evidence built in order, from bench accuracy to
prospective validation on the site's own population.

[DRAFTING INSTRUCTIONS: in the full-framework, build the competence case from the
author's assurance lineage in \texttt{review/references/author\_works.bib}
(\texttt{Kawchak2026MobilePancreaticCancerUnitree},
\texttt{Kawchak2026VVUQOncologyClinicalTrial}) and the 172 of 172 automated-test
result carried in the prior review. Anchor the human-plus-machine framing in
\texttt{topol2019high}. Reproduce Mermaid
\texttt{adoption/mermaid/diagrams/06-competence-evidence-stack.md} and the
calibrated-trust band \texttt{05-calibrated-trust-quadrant.md} as colored cfwfig
TikZ. Carry the ADOPTION ACTION ``local validation''. Build a body-width table of
evidence levels (bench, retrospective, prospective, assurance) with what each
establishes.]

\begin{cfwfig}
\node[cfone] (a) at (0,0) {Bench}; \node[cftwo] (b) at (3.0,0) {Retrospective};
\node[cftwo] (c) at (6.4,0) {Prospective}; \node[cfact] (d) at (9.8,0) {Assurance 172/172};
\node[cfhope] (e) at (6.4,-2.0) {Competence established};
\draw[cfedge] (a)--(b); \draw[cfedge] (b)--(c); \draw[cfedge] (c)--(d);
\draw[cfedge] (d) to[out=-90,in=0] (e.east);
\end{cfwfig}
\vspace{-0.2cm}
\cfwcaption{Figure 3. The competence evidence stack (placeholder; expanded from
Mermaid 06 in the full-framework).}

Competence is answered not by a claim but by evidence the clinician can inspect,
and by the discipline of validating on the local population before clinical use.
